# Template de saída para terminal do solver_richards2.py
# Salve como tela.tex e personalize conforme desejar

====  Iniciando simulação de drenagem 2D (Richards)  ====

Arquivo de malha lido: malha_dominio.xdmf
     - Pontos (nós): {num_nodes}
     - Elementos (triângulos): {num_elems}

Parâmetros do problema:
     - Condutividade hidráulica saturada (Ksat): {ksat:.2e} m/s
     - Gravidade (g): {gravidade:.2f} m/s²
     - Passo de tempo (dt): {dt:.1f} s
     - Tempo total (T): {tempo_total:.1f} s

Condições de contorno:
     - Topo (y = max): h = {h_top:.2f} m
     - Base (y = min): h = {h_bottom:.2f} m

Executando solver FEM 2D...
     - Sistema linear montado.
     - Solução obtida.

Exportando resultados:
     - Arquivo ParaView: saida_solver2.vtu

Resumo do campo de pressão exportado:
     - h (min, max, mean): {h_min:.2f}, {h_max:.2f}, {h_mean:.2f}

====  Simulação finalizada com sucesso  ====
